% TOP_PROBLEM: {"id": 1100206, "title": "Questions in Geometric Group Theory — Q 2.6", "category_id": 17, "category": {"id": 17, "name": "group_theory", "display_name": "Group Theory", "description": "Problems about groups, group actions, representations, and related algebraic structures.", "slug": "group-theory", "order_index": 17, "created_at": "2026-07-31T15:26:25.671Z"}, "status": "open", "problem_number": "AMR-010-0206", "set_id": 13, "statusQualification": "States the standard proper CAT(0) setting in which visual boundaries are compact metric spaces; asks for the source common-domain version of cell-like equivalence."}
% FIRST_POSTED: 2026-10-01
% ENTRY_KIND: statement_only
% UnsolvedMath ID 1100206; source catalogue code AMR-010-0206
% !TeX program = lualatex
% Definition and sources only; this file is not an LLM solution attempt.
\documentclass[11pt,a4paper]{article}
\usepackage[margin=25mm]{geometry}
\usepackage{fontspec}
\setmainfont{lmroman10-regular.otf}[BoldFont=lmroman10-bold.otf,
 ItalicFont=lmroman10-italic.otf,BoldItalicFont=lmroman10-bolditalic.otf]
\usepackage{amsmath,amssymb,mathtools}
\usepackage{unicode-math}
\setmathfont{latinmodern-math.otf}
\AtBeginDocument{\let\setminus\smallsetminus}
\usepackage{xurl,hyperref}
\hypersetup{colorlinks=true,urlcolor=blue}
\setcounter{secnumdepth}{0}
\setlength{\parindent}{0pt}
\setlength{\parskip}{5pt}
\setlength{\emergencystretch}{3em}
\begin{document}
\section*{Questions in Geometric Group Theory — Q 2.6}\label{1100206}
{\small UnsolvedMath ID 1100206 · AMR-010-0206 · Group Theory · AMR Open Problem Lists}
\subsection{Definitions and mathematical statement}
Let a discrete group $G$ act geometrically on two proper CAT$(0)$ geodesic spaces $X$ and $Y$. Here proper means that closed bounded balls are compact, and geometrically means by isometries, properly discontinuously, and with compact quotient. The CAT$(0)$ condition compares geodesic triangles with Euclidean triangles, requiring distances between corresponding side points not to increase.

The visual boundary $\partial X$ consists of geodesic rays modulo finite Hausdorff distance. Based at a point of $X$, each class has a unique representative, and uniform convergence on bounded time intervals defines the cone topology. This is a compact metrizable boundary; define $\partial Y$ similarly.

A compact metric space $C$ is cell-like if, after embedding it in the Hilbert cube, its inclusion into every neighborhood is null-homotopic in that neighborhood. This is the shape-of-a-point condition, independent of the embedding. A cell-like map is a continuous surjection whose point preimages are cell-like compacta.

Must there be a compact metric space $Z$ and cell-like maps
\[
 Z\longrightarrow\partial X,\qquad Z\longrightarrow\partial Y?
\]
No action of $G$ on $Z$, and no equivariance of these maps, is required. The question asks for a common cell-like resolution, even when the two boundaries are not homeomorphic. Properness of the spaces makes explicit the usual compact visual-boundary setting; it is not a request to identify the spaces themselves by an isometry.
\subsection{Short English statement}
Do two geometric CAT(0) actions of the same group have visual boundaries admitting a common cell-like resolution?
\subsection{Sources}
\begin{itemize}
\item[S1] ULAM.AI and the UnsolvedMath Contributors, supplied problems.json, record 1100206 (AMR-010-0206), AMR Open Problem Lists. Catalogue identity and source transcription; CC BY 4.0.
\url{https://huggingface.co/datasets/ulamai/UnsolvedMath}.
\item[S2] Bestvina - Questions in Geometric Group Theory (pdf) (2004), Question 2.6 (PDF page 7). Original source indexed by AMR-010-0206.
\url{https://www.math.utah.edu/~bestvina/eprints/questions-updated.pdf}.
\end{itemize}
\end{document}
